\documentclass[Main.tex]{subfiles}
\begin{document}
\gdef\muctieu{%
	\begin{enumerate}
		\item Trình bày được tính chất và ứng dụng của một số vật liệu thông dụng.
		\item Trình bày được tính chất và ứng dụng của một số nguyên liệu thông dụng.
		\item Nêu được khái niệm nhiên liệu và phân loại được nhiên liệu theo trạng thái.
		\item Trình bày được tính chất và ứng dụng của một số nhiên liệu thông dụng.
		\item Nêu được cách sử dụng nhiên liệu an toàn, hiệu quả và trình bày sơ lược về an ninh năng lượng.
		\item Trình bày được vai trò của lương thực, thực phẩm và các nhóm chất dinh dưỡng chính.
		\item Nêu được cách bảo quản lương thực, thực phẩm và ảnh hưởng của chế độ dinh dưỡng tới sức khoẻ.
		\item Đề xuất được phương án sử dụng vật liệu, nguyên liệu, nhiên liệu, lương thực -- thực phẩm an toàn, hiệu quả, bảo đảm sự phát triển bền vững.
	\end{enumerate}
}%
\chapter{Một số vật liệu, nguyên liệu, nhiên liệu, lương thực – thực phẩm thông dụng}
\begin{Muctieu}
	\muctieu
\end{Muctieu}
\section{Một số nhiên liệu}
\subsection{Lý thuyết trọng tâm}
\begin{kienthuccannho}
	\subsubsection{Các loại nhiên liệu}
	\begin{itemize}
		\item \textbf{\textit{Nhiên liệu}} (chất đốt) là những chất cháy được và khi cháy tỏa nhiều nhiệt.
		\item Dựa vào trạng thái, nhiên liệu được chia thành ba loại:
	\end{itemize}
	\begin{center}
		\renewcommand{\arraystretch}{1.15}
		\begin{tabular}{|C{0.28\linewidth}|C{0.28\linewidth}|C{0.28\linewidth}|}
			\hline
			\textbf{Nhiên liệu rắn} & \textbf{Nhiên liệu lỏng} & \textbf{Nhiên liệu khí} \\
			\hline
			Than, gỗ & Xăng dầu, cồn & Gas, khí thiên nhiên \\
			\hline
		\end{tabular}
	\end{center}
	\subsubsection{Tính chất và ứng dụng của nhiên liệu}
	\begin{center}
		\renewcommand{\arraystretch}{1.15}
		\begin{tabular}{|C{0.14\linewidth}|L{0.42\linewidth}|L{0.28\linewidth}|}
			\hline
			\textbf{Nhiên liệu} & \textbf{Tính chất} & \textbf{Ứng dụng} \\
			\hline
			Than đá
			&
			\begin{itemize}
				\item Chất rắn, cháy tỏa nhiều nhiệt.
				\item Khi cháy sinh nhiều chất độc hại, đặc biệt là khí carbon monoxide (CO) $\Rightarrow$ Không sử dụng than trong phòng kín.
			\end{itemize}
			&
			\begin{itemize}
				\item Đun nấu, sưởi ấm, nhiên liệu trong công nghiệp.
			\end{itemize}
			\\
			\hline
			Xăng, dầu
			&
			\begin{itemize}
				\item Chất lỏng, cháy tỏa nhiều nhiệt.
				\item Từ dầu mỏ người ta tạo ra xăng dầu, khí hóa lỏng, …
			\end{itemize}
			&
			\begin{itemize}
				\item Đun nấu, chạy động cơ, …
			\end{itemize}
			\\
			\hline
			Khí gas
			&
			\begin{itemize}
				\item Chất khí, cháy tỏa nhiều nhiệt.
			\end{itemize}
			&
			\begin{itemize}
				\item Dùng trong đun nấu, …
			\end{itemize}
			\\
			\hline
		\end{tabular}
	\end{center}
	\subsubsection{Sử dụng nhiên liệu an toàn và hiệu quả}
	\begin{center}
		\renewcommand{\arraystretch}{1.15}
		\begin{tabular}{|L{0.43\linewidth}|L{0.43\linewidth}|}
			\hline
			\textbf{Lợi ích sử dụng nhiên liệu an toàn, hiệu quả}
			&
			\textbf{Biện pháp sử dụng nhiên liệu an toàn, hiệu quả}
			\\
			\hline
			\begin{itemize}
				\item Tránh cháy nổ gây nguy hiểm đến con người và tài sản.
				\item Giảm thiểu ô nhiễm môi trường.
				\item Làm nhiên liệu cháy hoàn toàn và tận dụng hiệu quả lượng nhiệt do quá trình cháy tạo ra.
			\end{itemize}
			&
			\begin{itemize}
				\item Cung cấp đủ oxygen cho quá trình cháy.
				\item Tăng diện tích tiếp xúc giữa không khí và nhiên liệu.
				\item Điều chỉnh nhiên liệu để duy trì sự cháy ở mức độ cần thiết, phù hợp với nhu cầu sử dụng.
			\end{itemize}
			\\
			\hline
		\end{tabular}
	\end{center}
	\subsubsection{Sơ lược về an ninh năng lượng}
	\begin{itemize}
		\item \textbf{\textit{An ninh năng lượng}} là sự đảm bảo đầy đủ năng lượng dưới nhiều dạng khác nhau, ưu tiên các nguồn năng lượng sạch và giá thành rẻ.
		\item \textbf{Giải pháp:} Sử dụng các nhiên liệu tái tạo như nhiên liệu sinh học, nhiên liệu xanh, thủy điện, năng lượng mặt trời, năng lượng gió, … thay thế cho nguồn năng lượng không tái tạo (than đá, dầu mỏ, khí thiên nhiên).
	\end{itemize}
\end{kienthuccannho}
\subsection{Bài tập}

%%%=================Phần bài tập tự luận===================%%%
\phan{Bài tập tự luận}
\textit{\large Thí sinh trả lời từ câu 1 đến câu 16}
\Opensolutionfile{ansbth}[Ans/LGBT-C03_B08_MOT_SO_NHIEN_LIEU]
\Opensolutionfile{ansbt}[Ans/AnsBT-C03_B08_MOT_SO_NHIEN_LIEU]

%%%=============BT_1=============%%%
\begin{bt}%[6K3N3-1]
	Em hãy kể tên các nhiên liệu được dùng trong các hình dưới đây.
	\begin{center}
		\renewcommand{\arraystretch}{1}
		\begin{tabular}{C{0.21\linewidth}C{0.21\linewidth}C{0.21\linewidth}C{0.21\linewidth}}
			\includegraphics[width=0.95\linewidth]{Images/DE_CUONG_HOA_6_HKI/nhien_lieu_c1_a.jpg}
			&
			\includegraphics[width=0.95\linewidth]{Images/DE_CUONG_HOA_6_HKI/nhien_lieu_c1_b.jpg}
			&
			\includegraphics[width=0.95\linewidth]{Images/DE_CUONG_HOA_6_HKI/nhien_lieu_c1_c.jpg}
			&
			\includegraphics[width=0.95\linewidth]{Images/DE_CUONG_HOA_6_HKI/nhien_lieu_c1_d.jpg}
			\\
			(a) & (b) & (c) & (d)
		\end{tabular}
	\end{center}
	\loigiai{
		Nhiên liệu đã dùng lần lượt là:
		\begin{itemize}
			\item (a) xăng (dầu diesel);
			\item (b) khí gas;
			\item (c) dầu hỏa;
			\item (d) gỗ.
		\end{itemize}
	}
\end{bt}

%%%=============BT_2=============%%%
\begin{bt}%[6K3H3-2]
	Hãy cho biết ứng dụng của các nhiên liệu: dầu hỏa, gỗ (gỗ vụn, mùn cưa, cành cây khô), xăng, than đá, khí thiên nhiên.
	\loigiai{
		\begin{center}
			\renewcommand{\arraystretch}{1.15}
			\begin{tabular}{|L{0.3\linewidth}|L{0.55\linewidth}|}
				\hline
				\textbf{Nhiên liệu} & \textbf{Ứng dụng} \\
				\hline
				Dầu hỏa & Đèn dầu, bếp dầu, động cơ xe, máy phát điện, … \\
				\hline
				Gỗ (gỗ vụn, mùn cưa, cành cây khô) & Làm củi đun nấu, sưởi ấm, … \\
				\hline
				Xăng & Chạy xe ô tô, máy phát điện, … \\
				\hline
				Than đá & Lò cao nung vôi, sản xuất xi măng, luyện gang, thép. \\
				\hline
				Khí thiên nhiên & Nấu ăn, chạy máy phát điện, lò nung gạch, gốm, … \\
				\hline
			\end{tabular}
		\end{center}
	}
\end{bt}

%%%=============BT_3=============%%%
\begin{bt}%[6K3H3-1]
	Em hãy cho biết các biểu tượng trong hình dưới đây chỉ loại nhiên liệu nào?
	\begin{center}
		\includegraphics[width=0.9\linewidth]{Images/DE_CUONG_HOA_6_HKI/bieu_tuong_nhien_lieu.png}
		\\
		(a) \qquad (b) \qquad (c) \qquad (d) \qquad (e) \qquad (g)
	\end{center}
	\loigiai{
		\begin{itemize}
			\item (a) Than.
			\item (b) Xăng dầu.
			\item (c) Gas.
			\item (d) Năng lượng gió.
			\item (e) Năng lượng mặt trời.
			\item (g) Khí thiên nhiên.
		\end{itemize}
	}
\end{bt}

%%%=============BT_4=============%%%
\begin{bt}%[6K3V3-3]
	Hãy kể tên các nhiên liệu thường dùng trong việc đun nấu và nêu cách dùng nhiên liệu đó an toàn, tiết kiệm.
	\loigiai{
		\begin{itemize}
			\item \textbf{Củi:} phơi khô, chẻ nhỏ củi khi đun nấu, sử dụng quạt gió khi nhóm củi.
			\item \textbf{Gas:} sử dụng bình gas chính hãng, có van an toàn, sử dụng nhỏ lửa khi thực phẩm bắt đầu chín, vệ sinh bếp gas thường xuyên.
		\end{itemize}
	}
\end{bt}

%%%=============BT_5=============%%%
\begin{bt}%[6K3V3-3]
	Các việc làm dưới đây có thể có nhược điểm hoặc tác hại gì?
	\begin{enumerate}
		\item[(a)] Đun nấu để ngọn lửa quá to, không phù hợp với mục đích sử dụng.
		\item[(b)] Đun bếp than trong phòng kín.
	\end{enumerate}
	\loigiai{
		\begin{itemize}
			\item (a) Đun nấu để ngọn lửa quá to và không phù hợp với mục đích sử dụng sẽ dễ gây lãng phí nhiên liệu, đồng thời mất an toàn cháy nổ.
			\item (b) Trong phòng kín, không khí khó lưu thông với bên ngoài, thậm chí không thể lưu thông với bên ngoài. Khi đó, việc đốt than làm lượng oxygen giảm và sinh ra khí độc là carbon monoxide, có thể gây ngạt, thậm chí gây tử vong.
		\end{itemize}
	}
\end{bt}

%%%=============BT_6=============%%%
\begin{bt}%[6K3V3-3]
	Giải thích tác dụng của các việc làm sau đây:
	\begin{enumerate}
		\item[(a)] Chẻ nhỏ củi khi đun nấu.
		\item[(b)] Tạo các lỗ trong viên than tổ ong.
		\item[(c)] Quạt gió vào bếp lò khi nhóm lửa.
		\item[(d)] Đậy bớt cửa lò khi ủ bếp.
	\end{enumerate}
	\loigiai{
		\begin{itemize}
			\item (a) Tăng diện tích tiếp xúc giữa củi và oxygen (trong không khí) làm cho củi dễ cháy.
			\item (b) Không khí dễ dàng chui vào các lỗ hổng của than để tăng diện tích tiếp xúc giữa than và oxygen làm cho than dễ cháy.
			\item (c) Quạt gió (không khí) vào bếp lò để bổ sung oxygen làm cho củi, than dễ cháy.
			\item (d) Khi lò nóng rồi, người ta đậy bớt cửa lò để không cho không khí vào nhiều, hạn chế cháy hết củi hoặc than, làm cho bếp giữ nóng được lâu.
		\end{itemize}
	}
\end{bt}

%%%=============BT_7=============%%%
\begin{bt}%[6K3V3-3]
	Gas là một chất rất dễ cháy, khi gas trộn lẫn oxygen trong không khí nó sẽ trở thành một hỗn hợp dễ nổ. Hỗn hợp này sẽ bốc cháy và nổ rất mạnh khi có tia lửa điện hoặc đánh lửa từ bật gas, bếp gas.
	\begin{enumerate}
		\item[(a)] Chúng ta nên làm gì sau khi sử dụng bếp gas để đảm bảo an toàn?
		\item[(b)] Tại sao nên để bình gas ở nơi thoáng khí?
		\item[(c)] Trong trường hợp đang nấu ăn mà vòi dẫn gas bị hở và gas phun ra, cháy mạnh thì ta nên làm thế nào?
		\item[(d)] Khi ta đi học về, mở cửa nhà ra mà ngửi thấy mùi gas thì em nên làm gì?
	\end{enumerate}
	\loigiai{
		\begin{itemize}
			\item (a) Sau khi sử dụng bếp gas thì nên khóa van an toàn để tránh trường hợp gas bị rò ra ngoài có thể gây cháy nổ.
			\item (b) Để bình gas nơi thoáng khí để khi lỡ có rò gas thì khí cũng bay ra xa, làm loãng lượng gas trong không gian nhà bếp và tránh được nguy cơ cháy nổ.
			\item (c) Khi vòi dẫn gas bị hở và cháy, cần bình tĩnh tránh xa ngọn lửa, sau đó vặn khóa van an toàn bình gas lại. Trong trường hợp ngọn lửa lớn không tiếp xúc được với khóa gas thì dùng chăn ướt tấp kín để dập tắt ngọn lửa rồi khóa van an toàn bình gas.
			\item (d) Đi học về mà ngửi thấy mùi gas thì nên hành động như sau:
			\begin{itemize}
				\item Mở hết cửa để gas bay ra ngoài.
				\item Khóa van an toàn ở bình gas.
				\item Tuyệt đối không bật công tắc điện, không đánh lửa.
				\item Báo cho người lớn để kiểm tra và sửa chữa trước khi sử dụng lại.
			\end{itemize}
		\end{itemize}
	}
\end{bt}

%%%=============BT_8=============%%%
\begin{bt}%[6K3H3-6]
	Tại sao phải sử dụng các nhiên liệu tái tạo thay thế dần cho các nguồn nhiên liệu hóa thạch?
	\loigiai{
		\begin{itemize}
			\item Nhiên liệu hóa thạch (than đá, khí tự nhiên, …) có trong lòng đất là có hạn, phải mất hàng trăm triệu năm mới bổ sung được, do đó nếu khai thác liên tục nhiên liệu hóa thạch sẽ làm cạn kiệt nguồn nhiên liệu.
			\item Nhiên liệu hóa thạch chứa hàm lượng lớn carbon nên khi cháy tạo ra khí carbon dioxide gây hiệu ứng nhà kính (làm Trái Đất nóng lên và gây biến đổi khí hậu) và khí độc carbon monoxide ảnh hưởng đến sức khỏe con người.
		\end{itemize}
		Vì vậy cần thay thế dần bằng các nhiên liệu tái tạo.
	}
\end{bt}

%%%=============BT_9=============%%%
\begin{bt}%[6K3C3-6]
	Đọc đoạn thông tin sau và trả lời các câu hỏi.
	\begin{center}
		\textbf{MỘT SỐ LOẠI NHIÊN LIỆU CỦA TƯƠNG LAI}
	\end{center}
	Trong những năm tới, rất có thể bạn sẽ thường xuyên thấy những chiếc ô tô chạy bằng những loại nhiên liệu dưới đây.

	\textbf{\textit{Hydrogen.}} Các nhà sản xuất đang lên kế hoạch nạp hydrogen vào ô tô như các loại xăng dầu thông thường. Khi đó hydrogen sẽ chuyển hóa năng lượng hóa học thành điện và cung cấp cho hoạt động của chiếc xe. Tất cả những gì xe thải ra trong quá trình vận hành sẽ chỉ là nước.

	\textbf{\textit{Dầu diesel sinh học.}} Diesel sinh học là loại nhiên liệu được sản xuất từ dầu thực vật hay mỡ động vật để trở thành nhiên liệu cho xe. Nó được đánh giá là một nhiên liệu sạch với mức khí thải thấp hơn nhiều so với các loại nhiên liệu thông thường. Hơn nữa, vì được sản xuất từ các nguyên liệu rẻ, sẵn có như đậu tương nên diesel sinh học giúp các quốc gia giảm sự phụ thuộc vào nguồn dầu nhập khẩu.

	\textbf{\textit{Nhiên liệu pha ethanol.}} Thông thường, ethanol được sản xuất từ quá trình lên men của ngũ cốc như ngô. Đây là một nguồn nhiên liệu sạch và sản sinh khí nhà kính thấp hơn so với các loại khác. Ethanol được đưa vào xe sau khi đã pha trộn với xăng tùy theo từng nồng độ khác nhau. Nhiều quốc gia hiện nay đang sử dụng E85 với tỉ lệ pha trộn $85\%$ ethanol và $15\%$ xăng về thể tích.
	\begin{enumerate}
		\item[(a)] Vì sao hydrogen được coi là nhiên liệu không gây ô nhiễm môi trường?
		\item[(b)] Sử dụng các nhiên liệu như hydrogen, dầu diesel sinh học, … có lợi gì đối với an ninh năng lượng của mỗi quốc gia?
		\item[(c)] Xăng E90 có tỉ lệ $90\%$ ethanol và $10\%$ xăng về thể tích. Người ta phải thêm bao nhiêu lít ethanol vào $1$ lít xăng E85 để có xăng E90? (Giả sử không có hao hụt thể tích khi pha trộn.)
	\end{enumerate}
	\loigiai{
		\begin{itemize}
			\item (a) Xe chạy bằng nhiên liệu hydrogen chỉ thải ra nước, không gây ô nhiễm môi trường.
			\item (b) Các quốc gia sẽ có những nguồn năng lượng sạch, rẻ, bảo đảm nhu cầu sử dụng, giảm sự phụ thuộc vào dầu nhập khẩu.
			\item (c) Trong $1$ lít xăng E85 có $0{,}15$ lít xăng và $0{,}85$ lít ethanol.

			Thêm ethanol vào $1$ lít xăng E85 để có xăng E90 thì thể tích xăng không thay đổi và thể tích ethanol gấp $9$ lần thể tích xăng.

			Thể tích ethanol trong xăng E90 (sau khi được pha chế) là $0{,}15 \times 9 = 1{,}35$ (lít).

			Thể tích ethanol thêm vào là $1{,}35 - 0{,}85 = 0{,}5$ (lít).
		\end{itemize}
		\textbf{\textit{Đáp số:} thêm $0{,}5$ lít ethanol.}
	}
\end{bt}

%%%=============BT_10=============%%%
\begin{bt}%[6K3N3-1]
	Em hãy cho biết nhiên liệu có thể tồn tại ở những thể nào, lấy ví dụ minh họa.
	\loigiai{
		Nhiên liệu có thể tồn tại ở ba thể:
		\begin{itemize}
			\item Thể rắn: than đá, gỗ.
			\item Thể lỏng: xăng, dầu.
			\item Thể khí: khí gas.
		\end{itemize}
	}
\end{bt}

%%%=============BT_11=============%%%
\begin{bt}%[6K3V3-3]
	Nêu ba việc nên làm và ba việc nên tránh để sử dụng các nhiên liệu an toàn, hiệu quả, phòng tránh nguy cơ cháy nổ ở gia đình em.
	\loigiai{
		Một số việc nên làm:
		\begin{itemize}
			\item Sau khi sử dụng xong phải khóa bình gas.
			\item Sử dụng xong bếp cần tắt bếp.
			\item Không để các nhiên liệu gần nguồn điện.
		\end{itemize}
		Một số việc nên tránh:
		\begin{itemize}
			\item Mở các thiết bị lò sưởi, máy sấy không đúng với nhu cầu sử dụng.
			\item Sử dụng lửa quá to và không đúng mục đích khi đun nấu.
			\item Đun bếp than ở nơi không khí khó lưu thông.
		\end{itemize}
	}
\end{bt}

%%%=============BT_12=============%%%
\begin{bt}%[6K3V3-2]
	Tại sao khi gió thổi mạnh vào đống lửa to thì nó càng cháy mạnh còn thổi vào ngọn nến thì nó tắt ngay?
	\loigiai{
		\begin{itemize}
			\item Khi thổi vào đống lửa to, gió cung cấp thêm nhiều oxygen nên đống lửa sẽ càng cháy mạnh hơn.
			\item Còn khi gió thổi mạnh vào ngọn nến, nó làm nhiệt độ ngọn nến hạ đột ngột xuống dưới nhiệt độ cháy nên ngọn nến sẽ tắt.
		\end{itemize}
	}
\end{bt}

%%%=============BT_13=============%%%
\begin{bt}%[6K3V3-3]
	Em hãy tìm hiểu và nêu cách sử dụng khí gas, xăng trong sinh hoạt gia đình (để đun nấu, nhiên liệu chạy xe máy, ô tô, …) an toàn, tiết kiệm.
	\loigiai{
		Nguyên tắc sử dụng nhiên liệu an toàn là nắm vững tính chất đặc trưng của từng nhiên liệu. Dùng đủ, đúng cách là cách tiết kiệm nhiên liệu.
		\\
		Ví dụ:
		\begin{itemize}
			\item Khi dùng than, củi hoặc gas để nấu ăn chỉ cần để lửa ở mức phù hợp với việc đun nấu, không để quá to hoặc quá lâu khi không cần thiết.
			\item Với những đoạn đường không quá xa nên đi bộ hoặc đi xe đạp để tiết kiệm nhiên liệu và tăng cường vận động, tốt cho sức khỏe.
			\item Hạn chế dùng các phương tiện cá nhân, tăng sử dụng phương tiện giao thông công cộng.
		\end{itemize}
	}
\end{bt}

%%%=============BT_14=============%%%
\begin{bt}%[6K3C3-4]
	Bạn Linh lấy $2$ chiếc đèn trong phòng thí nghiệm rồi cho dầu hỏa vào đèn $1$, cồn ethanol vào đèn $2$. Dùng bật lửa gas thắp cả $2$ đèn lên rồi lấy hai tấm kính trắng che phía trên ngọn lửa của $2$ đèn. Kết quả bạn thấy tấm kính trên ngọn lửa đèn dầu bị đen (có muội than), còn tấm trên ngọn lửa đèn cồn thì không bị đen.
	\begin{enumerate}
		\item[(a)] Tại sao phòng thí nghiệm chỉ sử dụng đèn cồn mà không sử dụng đèn dầu hỏa?
		\item[(b)] Tại sao tấm kính che trên ngọn đèn dầu bị đen còn tấm che trên ngọn đèn cồn không bị đen?
		\item[(c)] Tại sao khi thắp đèn dầu mà ta vặn bấc càng cao lên thì trên chụp đèn càng nhanh đen?
	\end{enumerate}
	\loigiai{
		\begin{itemize}
			\item (a) Trong phòng thí nghiệm sử dụng đèn cồn sẽ không có muội than, không làm đen ống nghiệm nên dễ quan sát hiện tượng thí nghiệm. Nếu sử dụng đèn dầu sẽ sinh ra muội than, làm đen ống nghiệm dẫn đến khó quan sát hiện tượng thí nghiệm.
			\item (b) Do thiếu oxygen nên dầu cháy không hoàn toàn (cần nhiều oxygen hơn ethanol) sinh ra muội than (carbon). Còn ethanol cháy hết, không có muội than.
			\item (c) Khi vặn bấc càng cao thì dầu lên theo bấc càng nhiều, oxygen càng thiếu nên muội than sinh ra càng nhiều, chụp đèn sẽ đen hơn.
		\end{itemize}
	}
\end{bt}

%%%=============BT_15=============%%%
\begin{bt}%[6K3V3-6]
	Ở nhiều vùng nông thôn, người ta xây dựng hầm biogas để thu gom chất thải động vật. Chất thải được gom vào hầm sẽ phân hủy, theo thời gian tạo ra biogas. Biogas chủ yếu là khí methane, ngoài ra còn một lượng nhỏ các khí như ammonia, hydrogen sulfide, sulfur dioxide, … Biogas tạo ra sẽ được thu lại và dẫn lên để làm nhiên liệu khí phục vụ cho đun nấu hoặc chạy máy phát điện.
	\begin{enumerate}
		\item[(a)] Theo em, việc xây hầm thu chất thải sản xuất biogas đem lại những lợi ích gì?
		\item[(b)] Nếu sử dụng trực tiếp, biogas thường sẽ có mùi hôi của các khí như ammonia, hydrogen sulfide, … Em hãy đề xuất biện pháp giảm thiểu mùi hôi đó.
	\end{enumerate}
	\loigiai{
		\begin{itemize}
			\item (a) Việc thu gom chất thải tạo biogas có nhiều tác dụng:
			\begin{itemize}
				\item Làm sạch môi trường, hạn chế gây ô nhiễm môi trường.
				\item Tiêu diệt mầm bệnh gây hại. Nếu chất thải động vật thải trực tiếp ra môi trường sẽ phát tán nhiều mầm bệnh.
				\item Thu được biogas làm nhiên liệu phục vụ cuộc sống, tiết kiệm tiền mua nhiên liệu.
			\end{itemize}
			\item (b) Để hạn chế mùi hôi cần loại bỏ một số khí có mùi hôi trong thành phần của biogas. Muốn vậy, ta có thể dẫn khí qua thùng chứa than hoạt tính để khử mùi trước khi đưa vào sử dụng, hoặc áp dụng quy trình sản xuất và thu biogas sạch.
		\end{itemize}
	}
\end{bt}

%%%=============BT_16=============%%%
\begin{bt}%[6K3V3-3]
	Em hãy tìm hiểu và thảo luận về các nguồn nhiên liệu hóa thạch của Việt Nam.
	\loigiai{
		\begin{itemize}
			\item Nhiên liệu hóa thạch được tạo thành bởi quá trình phân hủy của các xác động thực vật bị chôn vùi hàng triệu năm. Tùy thuộc môi trường và điều kiện phân hủy mà nhiên liệu hóa thạch hình thành dưới dạng than (dạng rắn), dầu (dạng lỏng) và khí thiên nhiên (dạng khí).
			\item Ở nước ta, dầu mỏ và khí thiên nhiên tập trung chủ yếu ở thềm lục địa phía nam như: mỏ dầu Bạch Hổ, Đại Hùng, Rồng, Rạng Đông, Lan Tây. Mỏ khí thiên nhiên được khai thác ở mỏ Tiền Hải, tỉnh Thái Bình. Mỏ than ở Quảng Ninh.
			\item Tốc độ khai thác, tiêu thụ nhiên liệu hóa thạch nhanh hơn rất nhiều so với thời gian hình thành nên cần phải sử dụng tiết kiệm nguồn nhiên liệu này.
		\end{itemize}
	}
\end{bt}
\Closesolutionfile{ansbt}
\Closesolutionfile{ansbth}

%%%=================Phần trắc nghiệm nhiều lựa chọn===================%%%
\phan{Bài tập trắc nghiệm nhiều lựa chọn}
\textit{\large Thí sinh trả lời từ câu 1 đến câu 25. Mỗi câu thí sinh chỉ chọn một phương án}
\Opensolutionfile{ansex}[Ans/LGEX-C03_B08_MOT_SO_NHIEN_LIEU]
\Opensolutionfile{ans}[Ans/Ans-C03_B08_MOT_SO_NHIEN_LIEU]

%%%=============EX_1=============%%%
\begin{ex}%[6K3N3-1]
	Nhiên liệu là
	\choice
	{một số chất hoặc hỗn hợp chất được dùng làm nguyên liệu đầu vào cho các quá trình sản xuất hoặc chế tạo}
	{những chất được oxi hóa để cung cấp năng lượng cho hoạt động của cơ thể sống}
	{những vật liệu dùng trong quá trình xây dựng}
	{\True những chất cháy được để cung cấp năng lượng dạng nhiệt hoặc ánh sáng nhằm phục vụ mục đích sử dụng của con người}
	\loigiai{Nhiên liệu (chất đốt) là những chất cháy được và khi cháy tỏa nhiều nhiệt, cung cấp năng lượng dạng nhiệt hoặc ánh sáng phục vụ con người.}
\end{ex}

%%%=============EX_2=============%%%
\begin{ex}%[6K3N3-1]
	Loại nhiên liệu nào dưới đây là nhiên liệu rắn?
	\choice
	{\True Than đá}
	{Dầu hỏa}
	{Dầu diesel}
	{Xăng}
	\loigiai{Than đá là chất rắn nên thuộc nhóm nhiên liệu rắn. Dầu hỏa, dầu diesel, xăng đều là chất lỏng.}
\end{ex}

%%%=============EX_3=============%%%
\begin{ex}%[6K3N3-1]
	Loại nhiên liệu nào dưới đây là nhiên liệu lỏng?
	\choice
	{Than đá}
	{\True Dầu hỏa}
	{Nến}
	{Gas}
	\loigiai{Dầu hỏa là chất lỏng nên thuộc nhóm nhiên liệu lỏng.}
\end{ex}

%%%=============EX_4=============%%%
\begin{ex}%[6K3N3-1]
	Loại nhiên liệu nào dưới đây là nhiên liệu khí?
	\choice
	{Xăng}
	{Dầu hỏa}
	{Nến}
	{\True Gas}
	\loigiai{Gas là chất khí nên thuộc nhóm nhiên liệu khí.}
\end{ex}

%%%=============EX_5=============%%%
\begin{ex}%[6K3N3-2]
	Tính chất nào dưới đây là tính chất chung của nhiên liệu?
	\choice
	{Nhẹ hơn nước}
	{Tan trong nước}
	{\True Cháy được}
	{Là chất rắn}
	\loigiai{Điểm chung của mọi nhiên liệu là cháy được và khi cháy tỏa nhiều nhiệt.}
\end{ex}

%%%=============EX_6=============%%%
\begin{ex}%[6K3H3-3]
	Để sử dụng nhiên liệu tiết kiệm và hiệu quả cần phải cung cấp một lượng không khí hoặc oxygen
	\choice
	{\True vừa đủ}
	{thiếu}
	{dư}
	{tùy ý}
	\loigiai{Cung cấp lượng oxygen vừa đủ giúp nhiên liệu cháy hoàn toàn, tránh lãng phí và hạn chế sinh khí độc.}
\end{ex}

%%%=============EX_7=============%%%
\begin{ex}%[6K3N3-1]
	Nhiên liệu hóa thạch
	\choice
	{là nguồn nhiên liệu tái tạo}
	{là đá chứa ít nhất $50\%$ xác động và thực vật}
	{chỉ bao gồm dầu mỏ, than đá}
	{\True là nhiên liệu hình thành từ xác sinh vật bị chôn vùi và biến đổi hàng triệu năm trước}
	\loigiai{Nhiên liệu hóa thạch hình thành từ xác động thực vật bị chôn vùi và biến đổi qua hàng triệu năm, gồm than đá, dầu mỏ và khí thiên nhiên.}
\end{ex}

%%%=============EX_8=============%%%
\begin{ex}%[6K3N3-1]
	Nhiên liệu nào sau đây là nhiên liệu hóa thạch?
	\choice
	{Gỗ}
	{\True Dầu mỏ}
	{Hydrogen}
	{Ethanol (cồn)}
	\loigiai{Dầu mỏ là nhiên liệu hóa thạch. Gỗ, hydrogen và ethanol thuộc nhóm nhiên liệu tái tạo.}
\end{ex}

%%%=============EX_9=============%%%
\begin{ex}%[6K3N3-3]
	An ninh năng lượng là
	\choice
	{sự đảm bảo đầy đủ năng lượng dưới dạng các nguồn năng lượng hóa thạch}
	{sự đảm bảo đầy đủ năng lượng dưới dạng các nguồn năng lượng sạch}
	{\True sự đảm bảo đầy đủ năng lượng dưới nhiều dạng khác nhau, ưu tiên các nguồn năng lượng sạch và giá thành rẻ}
	{sự đảm bảo đầy đủ năng lượng dưới nhiều dạng khác nhau, ưu tiên các nguồn năng lượng sẵn có và giá thành đắt}
	\loigiai{An ninh năng lượng là sự đảm bảo đầy đủ năng lượng dưới nhiều dạng khác nhau, ưu tiên nguồn năng lượng sạch và giá thành rẻ.}
\end{ex}

%%%=============EX_10=============%%%
\begin{ex}%[6K3N3-1]
	Nhiên liệu nào sau đây không phải nhiên liệu hóa thạch?
	\choice
	{Than đá}
	{Dầu mỏ}
	{Khí tự nhiên}
	{\True Ethanol (cồn)}
	\loigiai{Ethanol được sản xuất từ quá trình lên men ngũ cốc nên là nhiên liệu sinh học (tái tạo), không phải nhiên liệu hóa thạch.}
\end{ex}

%%%=============EX_11=============%%%
\begin{ex}%[6K3H3-3]
	Việc làm nào có thể bảo đảm an toàn khi sử dụng xăng?
	\choice
	{\True Vận chuyển xăng trong các thiết bị chuyên dụng}
	{Để xăng gần nguồn nhiệt}
	{Sử dụng điện thoại tại các trạm xăng}
	{Lưu trữ xăng trong các chai nhựa để tiện sử dụng}
	\loigiai{Xăng rất dễ cháy nên phải vận chuyển và lưu trữ trong thiết bị chuyên dụng, tránh xa nguồn nhiệt và tia lửa điện.}
\end{ex}

%%%=============EX_12=============%%%
\begin{ex}%[6K3H3-3]
	Để sử dụng nhiên liệu tiết kiệm và hiệu quả, cần điều chỉnh lượng gas khi đun nấu
	\choice
	{không thay đổi trong suốt quá trình sử dụng}
	{\True phù hợp với nhu cầu sử dụng}
	{luôn ở mức nhỏ nhất có thể}
	{luôn ở mức lớn nhất có thể}
	\loigiai{Điều chỉnh lượng gas phù hợp với nhu cầu giúp duy trì sự cháy ở mức cần thiết, vừa tiết kiệm vừa an toàn.}
\end{ex}

%%%=============EX_13=============%%%
\begin{ex}%[6K3H3-3]
	Để củi dễ cháy khi đun nấu, người ta không dùng biện pháp nào sau đây?
	\choice
	{Phơi củi cho thật khô}
	{Cung cấp đầy đủ oxygen cho quá trình cháy}
	{\True Xếp củi chồng lên nhau, càng sít nhau càng tốt}
	{Chẻ nhỏ củi}
	\loigiai{Xếp củi quá sít nhau làm giảm diện tích tiếp xúc với không khí, oxygen khó lọt vào nên củi khó cháy.}
\end{ex}

%%%=============EX_14=============%%%
\begin{ex}%[6K3H3-3]
	Để sử dụng gas tiết kiệm, hiệu quả người ta sử dụng biện pháp nào sau đây?
	\choice
	{\True Tùy nhiệt độ cần thiết để điều chỉnh lượng gas}
	{Tốt nhất nên để gas ở mức độ cao nhất}
	{Tốt nhất nên để gas ở mức độ nhỏ nhất}
	{Ngăn không cho khí gas tiếp xúc với carbon dioxide}
	\loigiai{Điều chỉnh lượng gas theo nhiệt độ cần thiết giúp duy trì sự cháy phù hợp với nhu cầu, tiết kiệm nhiên liệu.}
\end{ex}

%%%=============EX_15=============%%%
\begin{ex}%[6K3H3-6]
	Xăng sinh học E5 ra đời góp phần tích cực vào vấn đề bảo vệ môi trường và làm giảm tác động mà xăng xe thông thường gây ra. Thành phần của xăng sinh học E5 gồm $95\%$ xăng A92 thông thường pha trộn với $5\%$ chất nào sau đây?
	\choice
	{Nước}
	{Dầu mỏ}
	{Khí gas}
	{\True Ethanol}
	\loigiai{Xăng E5 gồm $95\%$ xăng A92 và $5\%$ ethanol về thể tích.}
\end{ex}

%%%=============EX_16=============%%%
\begin{ex}%[6K3V3-3]
	Cho các việc làm sau:
	\begin{enumerate}
		\item[(a)] Sử dụng bếp gas phải có van an toàn, dùng xong phải khóa van lại.
		\item[(b)] Sử dụng bếp than trong phòng kín để đảm bảo lượng nhiệt không bị thất thoát.
		\item[(c)] Khi đi xe máy gặp đèn đỏ thời gian $100$ s thì không được tắt máy.
		\item[(d)] Tạo các lỗ trong viên than tổ ong để than cháy tốt hơn.
		\item[(e)] Vặn nhỏ lửa bếp gas khi thức ăn gần chín.
	\end{enumerate}
	Số việc làm nhằm sử dụng nhiên liệu an toàn, tiết kiệm là
	\choice
	{$2$}
	{\True $3$}
	{$4$}
	{$5$}
	\loigiai{Các việc làm đúng là (a), (d), (e). Việc (b) sai vì đốt than trong phòng kín sinh khí độc carbon monoxide. Việc (c) sai vì dừng đèn đỏ lâu nên tắt máy để tiết kiệm nhiên liệu. Vậy có $3$ việc làm.}
\end{ex}

%%%=============EX_17=============%%%
\begin{ex}%[6K3N3-1]
	Chất nào sau đây không phải là nhiên liệu?
	\choice
	{Than đá}
	{Xăng}
	{\True Nước}
	{Khí gas}
	\loigiai{Nước không cháy được nên không phải là nhiên liệu.}
\end{ex}

%%%=============EX_18=============%%%
\begin{ex}%[6K3N3-6]
	Nhiên liệu nào sau đây thuộc loại nhiên liệu tái tạo?
	\choice
	{Than đá}
	{Dầu mỏ}
	{Khí thiên nhiên}
	{\True Nhiên liệu sinh học}
	\loigiai{Nhiên liệu sinh học được sản xuất từ thực vật, động vật nên có thể tái tạo. Than đá, dầu mỏ, khí thiên nhiên là nhiên liệu hóa thạch, không tái tạo.}
\end{ex}

%%%=============EX_19=============%%%
\begin{ex}%[6K3H3-3]
	Khi đun nấu, người ta chẻ nhỏ củi nhằm mục đích nào sau đây?
	\choice
	{\True Tăng diện tích tiếp xúc giữa củi và oxygen trong không khí}
	{Làm giảm nhiệt độ cháy của củi}
	{Làm cho củi nặng hơn}
	{Ngăn không cho củi tiếp xúc với không khí}
	\loigiai{Chẻ nhỏ củi làm tăng diện tích tiếp xúc giữa củi và oxygen nên củi cháy dễ và cháy hoàn toàn hơn.}
\end{ex}

%%%=============EX_20=============%%%
\begin{ex}%[6K3H3-3]
	Vì sao không nên đun bếp than trong phòng kín?
	\choice
	{Vì than cháy sẽ làm mát không khí trong phòng}
	{\True Vì lượng oxygen giảm và sinh ra khí độc carbon monoxide, có thể gây ngạt}
	{Vì than không cháy được trong phòng kín}
	{Vì than sẽ hút hết hơi nước trong phòng}
	\loigiai{Trong phòng kín, không khí khó lưu thông nên lượng oxygen giảm, than cháy không hoàn toàn sinh ra khí độc carbon monoxide gây ngạt, thậm chí tử vong.}
\end{ex}

%%%=============EX_21=============%%%
\begin{ex}%[6K3N3-2]
	Trong công nghiệp, than đá được dùng chủ yếu cho công việc nào sau đây?
	\choice
	{Làm dây dẫn điện}
	{\True Nung vôi, sản xuất xi măng, luyện gang thép}
	{Sản xuất đồ nhựa gia dụng}
	{Làm lốp xe, gioăng cao su}
	\loigiai{Than đá cháy tỏa nhiều nhiệt nên được dùng trong lò cao nung vôi, sản xuất xi măng, luyện gang, thép.}
\end{ex}

%%%=============EX_22=============%%%
\begin{ex}%[6K3N3-2]
	Xăng, dầu được tách ra từ nguồn nào sau đây?
	\choice
	{Than đá}
	{\True Dầu mỏ}
	{Quặng bauxite}
	{Đá vôi}
	\loigiai{Từ dầu mỏ, người ta tách ra xăng, dầu hỏa, dầu diesel, khí hóa lỏng, …}
\end{ex}

%%%=============EX_23=============%%%
\begin{ex}%[6K3H3-6]
	Biện pháp nào sau đây góp phần bảo đảm an ninh năng lượng?
	\choice
	{Tăng cường khai thác tối đa than đá và dầu mỏ}
	{\True Tăng cường sử dụng năng lượng mặt trời, năng lượng gió, nhiên liệu sinh học}
	{Hạn chế nghiên cứu các nguồn năng lượng mới}
	{Đun nấu bằng than tổ ong trong khu dân cư}
	\loigiai{An ninh năng lượng ưu tiên các nguồn năng lượng sạch, giá thành rẻ nên cần tăng cường sử dụng năng lượng tái tạo thay cho nhiên liệu hóa thạch.}
\end{ex}

%%%=============EX_24=============%%%
\begin{ex}%[6K3V3-3]
	Khi đi học về, mở cửa nhà ra mà ngửi thấy mùi gas, việc làm nào sau đây tuyệt đối KHÔNG được thực hiện?
	\choice
	{Mở hết các cửa cho gas bay ra ngoài}
	{\True Bật công tắc điện để kiểm tra nguồn rò rỉ}
	{Khóa van an toàn ở bình gas}
	{Báo cho người lớn kiểm tra và sửa chữa}
	\loigiai{Gas trộn với không khí tạo hỗn hợp dễ nổ, chỉ cần một tia lửa điện khi bật công tắc cũng có thể gây cháy nổ. Vì vậy tuyệt đối không bật công tắc điện, không đánh lửa.}
\end{ex}

%%%=============EX_25=============%%%
\begin{ex}%[6K3V3-1]
	Cho các nhiên liệu: than đá, gỗ, xăng, dầu hỏa, khí gas, khí thiên nhiên, cồn. Số nhiên liệu ở thể lỏng là
	\choice
	{$2$}
	{\True $3$}
	{$4$}
	{$5$}
	\loigiai{Nhiên liệu ở thể lỏng gồm xăng, dầu hỏa, cồn. Than đá và gỗ ở thể rắn; khí gas và khí thiên nhiên ở thể khí. Vậy có $3$ nhiên liệu ở thể lỏng.}
\end{ex}
\Closesolutionfile{ans}
\Closesolutionfile{ansex}
%\bangdapan{Ans-C03_B08_MOT_SO_NHIEN_LIEU}

%%%=================Phần trắc nghiệm đúng sai===================%%%
\phan{Bài tập trắc nghiệm đúng sai}
\textit{\large Thí sinh trả lời từ câu 1 đến câu 10. Trong mỗi ý a), b), c), d) ở mỗi câu thí sinh chọn đúng hoặc sai}
\Opensolutionfile{ansex}[Ans/LGTF-C03_B08_MOT_SO_NHIEN_LIEU]
\Opensolutionfile{ansbook}[Ansbook/AnsTF-C03_B08_MOT_SO_NHIEN_LIEU]
\Opensolutionfile{ans}[Ans/Tempt-C03_B08_MOT_SO_NHIEN_LIEU]
\setcounter{ex}{0}

%%%=============TF_1=============%%%
\begin{ex}%[6K3N3-1]
	Xét các phát biểu về nhiên liệu.
	\choiceTF
	{Nhiên liệu là những chất cháy được và khi cháy thu nhiều nhiệt}
	{\True Nhiên liệu rắn như than, gỗ}
	{Nhiên liệu lỏng như xăng dầu, cồn, nước}
	{\True Nhiên liệu khí như gas, khí thiên nhiên}
	\loigiai{
		\begin{itemchoice}[F1,T2,F3,T4]
			\itemch Sai vì nhiên liệu cháy tỏa nhiều nhiệt chứ không thu nhiệt.
			\itemch Than và gỗ đều ở thể rắn nên là nhiên liệu rắn.
			\itemch Sai vì nước không cháy được nên nước không phải nhiên liệu.
			\itemch Gas và khí thiên nhiên ở thể khí nên là nhiên liệu khí.
	\end{itemchoice}}
\end{ex}

%%%=============TF_2=============%%%
\begin{ex}%[6K3H3-2]
	Xét các phát biểu về tính chất và ứng dụng của than đá và xăng, dầu.
	\choiceTF
	{\True Than đá là chất rắn, cháy tỏa nhiều nhiệt, khi cháy sinh ra nhiều chất độc hại đặc biệt là khí CO}
	{\True Than đá dùng để đun nấu, sưởi ấm và làm nhiên liệu trong công nghiệp}
	{Xăng, dầu là chất lỏng, cháy tỏa nhiều nhiệt, được tách ra từ mỏ than}
	{\True Xăng, dầu dùng để chạy động cơ}
	\loigiai{
		\begin{itemchoice}[T1,T2,F3,T4]
			\itemch Than đá là chất rắn, cháy tỏa nhiều nhiệt và sinh ra khí độc carbon monoxide (CO).
			\itemch Than đá được dùng để đun nấu, sưởi ấm và làm nhiên liệu trong công nghiệp.
			\itemch Sai vì xăng dầu được tách ra từ dầu mỏ, không phải từ mỏ than.
			\itemch Xăng, dầu được dùng để chạy động cơ xe máy, ô tô, máy phát điện, …
	\end{itemchoice}}
\end{ex}

%%%=============TF_3=============%%%
\begin{ex}%[6K3H3-3]
	Xét các phát biểu về lợi ích sử dụng nhiên liệu an toàn, hiệu quả.
	\choiceTF
	{\True Tránh cháy nổ gây nguy hiểm đến con người và tài sản}
	{\True Giảm thiểu ô nhiễm môi trường}
	{Tăng khả năng miễn dịch của cơ thể}
	{\True Làm nhiên liệu cháy hoàn toàn và tận dụng hiệu quả lượng nhiệt do quá trình cháy tạo ra}
	\loigiai{
		\begin{itemchoice}[T1,T2,F3,T4]
			\itemch Sử dụng nhiên liệu an toàn giúp tránh cháy nổ gây nguy hiểm cho người và tài sản.
			\itemch Nhiên liệu cháy hoàn toàn sinh ít chất độc hại nên giảm ô nhiễm môi trường.
			\itemch Sai vì việc sử dụng nhiên liệu không liên quan đến khả năng miễn dịch của cơ thể.
			\itemch Sử dụng hiệu quả giúp nhiên liệu cháy hoàn toàn và tận dụng tốt lượng nhiệt tạo ra.
	\end{itemchoice}}
\end{ex}

%%%=============TF_4=============%%%
\begin{ex}%[6K3H3-3]
	Xét các phát biểu về biện pháp sử dụng nhiên liệu an toàn, hiệu quả.
	\choiceTF
	{\True Cần cung cấp đủ oxygen cho quá trình cháy}
	{\True Tăng diện tích tiếp xúc giữa không khí và nhiên liệu}
	{\True Điều chỉnh nhiên liệu để duy trì sự cháy ở mức độ cần thiết, phù hợp với nhu cầu sử dụng}
	{Cần tận dụng triệt để các nhiên liệu có sẵn trong tự nhiên như than, xăng, dầu, …}
	\loigiai{
		\begin{itemchoice}[T1,T2,T3,F4]
			\itemch Cung cấp đủ oxygen giúp nhiên liệu cháy hoàn toàn.
			\itemch Tăng diện tích tiếp xúc giữa không khí và nhiên liệu giúp nhiên liệu dễ cháy hơn.
			\itemch Duy trì sự cháy ở mức cần thiết giúp tiết kiệm nhiên liệu.
			\itemch Sai vì việc sử dụng triệt để các nhiên liệu hóa thạch sẽ gây cạn kiệt nguồn tài nguyên và gây ô nhiễm môi trường.
	\end{itemchoice}}
\end{ex}

%%%=============TF_5=============%%%
\begin{ex}%[6K3H3-6]
	Xét các phát biểu về an ninh năng lượng.
	\choiceTF
	{\True An ninh năng lượng là sự đảm bảo đầy đủ năng lượng dưới nhiều dạng khác nhau, ưu tiên các nguồn năng lượng sạch và giá thành rẻ}
	{Cần tăng cường sử dụng các nhiên liệu tái tạo như nhiên liệu sinh học, năng lượng mặt trời, dầu mỏ, …}
	{\True Hạn chế sử dụng các nguồn năng lượng không tái tạo như than đá, khí thiên nhiên}
	{\True Tăng cường sản xuất và sử dụng xe máy, ô tô chạy điện thay vì chạy xăng}
	\loigiai{
		\begin{itemchoice}[T1,F2,T3,T4]
			\itemch Đây chính là khái niệm an ninh năng lượng.
			\itemch Sai vì dầu mỏ là nguồn năng lượng không tái tạo.
			\itemch Than đá, khí thiên nhiên là nhiên liệu hóa thạch, cần hạn chế sử dụng.
			\itemch Xe chạy điện giúp giảm tiêu thụ xăng dầu và giảm ô nhiễm môi trường.
	\end{itemchoice}}
\end{ex}

%%%=============TF_6=============%%%
\begin{ex}%[6K3N3-1]
	Xét các phát biểu về phân loại nhiên liệu.
	\choiceTF
	{\True Than và gỗ thuộc nhóm nhiên liệu rắn}
	{\True Cồn và xăng dầu thuộc nhóm nhiên liệu lỏng}
	{\True Khí gas và khí thiên nhiên thuộc nhóm nhiên liệu khí}
	{Nước là nhiên liệu lỏng vì nước tồn tại ở thể lỏng}
	\loigiai{
		\begin{itemchoice}[T1,T2,T3,F4]
			\itemch Than và gỗ đều ở thể rắn và cháy được nên là nhiên liệu rắn.
			\itemch Cồn, xăng dầu ở thể lỏng và cháy được nên là nhiên liệu lỏng.
			\itemch Khí gas, khí thiên nhiên ở thể khí và cháy được nên là nhiên liệu khí.
			\itemch Sai vì nhiên liệu phải cháy được, mà nước không cháy được.
	\end{itemchoice}}
\end{ex}

%%%=============TF_7=============%%%
\begin{ex}%[6K3V3-3]
	Xét các phát biểu về sử dụng nhiên liệu an toàn trong gia đình.
	\choiceTF
	{\True Sau khi dùng bếp gas xong cần khóa van an toàn của bình gas}
	{Nên để bình gas ở nơi kín gió để tránh thất thoát khí gas}
	{\True Khi ngửi thấy mùi gas cần mở hết cửa cho thoáng và không bật công tắc điện}
	{Có thể lưu trữ xăng trong chai nhựa đặt ở bếp để tiện sử dụng}
	\loigiai{
		\begin{itemchoice}[T1,F2,T3,F4]
			\itemch Khóa van an toàn sau khi dùng giúp tránh rò gas gây cháy nổ.
			\itemch Sai vì cần để bình gas nơi thoáng khí để khi rò rỉ, gas được làm loãng và bay ra xa.
			\itemch Mở cửa cho thoáng và không đánh lửa là thao tác đúng khi phát hiện mùi gas.
			\itemch Sai vì xăng rất dễ cháy, phải lưu trữ trong thiết bị chuyên dụng và tránh xa nguồn nhiệt.
	\end{itemchoice}}
\end{ex}

%%%=============TF_8=============%%%
\begin{ex}%[6K3H3-1]
	Xét các phát biểu về nhiên liệu hóa thạch.
	\choiceTF
	{\True Nhiên liệu hóa thạch hình thành từ xác sinh vật bị chôn vùi và biến đổi hàng triệu năm}
	{\True Than đá, dầu mỏ, khí thiên nhiên đều là nhiên liệu hóa thạch}
	{Nhiên liệu hóa thạch là nguồn nhiên liệu vô tận}
	{\True Đốt nhiên liệu hóa thạch sinh ra khí carbon dioxide gây hiệu ứng nhà kính}
	\loigiai{
		\begin{itemchoice}[T1,T2,F3,T4]
			\itemch Nhiên liệu hóa thạch được tạo thành từ xác động thực vật bị chôn vùi hàng triệu năm.
			\itemch Ba loại nhiên liệu hóa thạch chính là than đá (rắn), dầu mỏ (lỏng), khí thiên nhiên (khí).
			\itemch Sai vì trữ lượng nhiên liệu hóa thạch có hạn, khai thác liên tục sẽ cạn kiệt.
			\itemch Nhiên liệu hóa thạch chứa nhiều carbon nên khi cháy tạo khí carbon dioxide gây hiệu ứng nhà kính.
	\end{itemchoice}}
\end{ex}

%%%=============TF_9=============%%%
\begin{ex}%[6K3H3-6]
	Xét các phát biểu về nhiên liệu tái tạo.
	\choiceTF
	{\True Nhiên liệu sinh học, năng lượng mặt trời, năng lượng gió là các nguồn năng lượng tái tạo}
	{\True Xăng E5 là hỗn hợp gồm $95\%$ xăng A92 và $5\%$ ethanol về thể tích}
	{Than đá là nhiên liệu tái tạo vì có thể khai thác liên tục}
	{\True Sử dụng nhiên liệu tái tạo góp phần bảo đảm an ninh năng lượng}
	\loigiai{
		\begin{itemchoice}[T1,T2,F3,T4]
			\itemch Đây đều là các nguồn năng lượng có thể tái tạo, thay thế cho nhiên liệu hóa thạch.
			\itemch Xăng E5 gồm $95\%$ xăng A92 và $5\%$ ethanol về thể tích.
			\itemch Sai vì than đá là nhiên liệu hóa thạch, cần hàng trăm triệu năm mới hình thành nên không tái tạo.
			\itemch Nhiên liệu tái tạo là nguồn năng lượng sạch, giá thành rẻ nên góp phần bảo đảm an ninh năng lượng.
	\end{itemchoice}}
\end{ex}

%%%=============TF_10=============%%%
\begin{ex}%[6K3H3-2]
	Xét các phát biểu về tính chất và ứng dụng của một số nhiên liệu.
	\choiceTF
	{\True Than đá là chất rắn, cháy tỏa nhiều nhiệt}
	{Mọi nhiên liệu sau khi cháy đều có thể tái sử dụng}
	{\True Dầu hỏa được dùng cho đèn dầu, bếp dầu, động cơ xe, máy phát điện}
	{Khí thiên nhiên là chất lỏng, được dùng để chạy máy phát điện}
	\loigiai{
		\begin{itemchoice}[T1,F2,T3,F4]
			\itemch Than đá ở thể rắn và cháy tỏa nhiều nhiệt.
			\itemch Sai vì nhiên liệu sau khi cháy đã biến đổi thành chất khác nên không tái sử dụng được.
			\itemch Dầu hỏa được dùng cho đèn dầu, bếp dầu, động cơ xe, máy phát điện.
			\itemch Sai vì khí thiên nhiên là chất khí, không phải chất lỏng.
	\end{itemchoice}}
\end{ex}
\Closesolutionfile{ans}
\Closesolutionfile{ansbook}
\Closesolutionfile{ansex}
%\bangdapanTF{AnsTF-C03_B08_MOT_SO_NHIEN_LIEU}

%%%=================Phần trắc nghiệm trả lời ngắn===================%%%
\phan{Bài tập trắc nghiệm trả lời ngắn}
\textit{\large Thí sinh trả lời từ câu 1 đến câu 10}
\Opensolutionfile{ansex}[Ans/LGSA-C03_B08_MOT_SO_NHIEN_LIEU]
\Opensolutionfile{ansexh}[Ans/AnsSA-C03_B08_MOT_SO_NHIEN_LIEU]
\setcounter{ex}{0}

%%%=============SA_1=============%%%
\begin{ex}%[6K3N3-1]
	Cho các chất: $(1)$ than, $(2)$ nước, $(3)$ cồn, $(4)$ đá vôi, $(5)$ muối ăn. Hãy sắp xếp các chất là nhiên liệu thành bộ số theo số thứ tự tăng dần (ví dụ $12$, $134$, …).
	\shortans{$13$}
	\loigiai{
		\begin{itemize}
			\item Than $(1)$ và cồn $(3)$ đều cháy được và khi cháy tỏa nhiều nhiệt nên là nhiên liệu.
			\item Nước, đá vôi, muối ăn không cháy được nên không phải nhiên liệu.
		\end{itemize}
		\textbf{\textit{Đáp số:} $13$.}
	}
\end{ex}

%%%=============SA_2=============%%%
\begin{ex}%[6K3N3-1]
	Cho các nhiên liệu: $(1)$ than, $(2)$ gỗ, $(3)$ xăng, $(4)$ gas, $(5)$ cồn. Hãy sắp xếp các nhiên liệu ở trạng thái rắn thành bộ số theo số thứ tự tăng dần (ví dụ $12$, $134$, …).
	\shortans{$12$}
	\loigiai{
		\begin{itemize}
			\item Than $(1)$ và gỗ $(2)$ ở trạng thái rắn.
			\item Xăng, cồn ở trạng thái lỏng; gas ở trạng thái khí.
		\end{itemize}
		\textbf{\textit{Đáp số:} $12$.}
	}
\end{ex}

%%%=============SA_3=============%%%
\begin{ex}%[6K3N3-1]
	Cho các nhiên liệu: $(1)$ than, $(2)$ khí thiên nhiên, $(3)$ gỗ, $(4)$ xăng, $(5)$ gas, $(6)$ cồn. Hãy sắp xếp các nhiên liệu ở trạng thái lỏng thành bộ số theo số thứ tự tăng dần (ví dụ $12$, $134$, …).
	\shortans{$46$}
	\loigiai{
		\begin{itemize}
			\item Xăng $(4)$ và cồn $(6)$ ở trạng thái lỏng.
			\item Than, gỗ ở trạng thái rắn; khí thiên nhiên, gas ở trạng thái khí.
		\end{itemize}
		\textbf{\textit{Đáp số:} $46$.}
	}
\end{ex}

%%%=============SA_4=============%%%
\begin{ex}%[6K3H3-2]
	Cho các phát biểu sau:
	\begin{enumerate}[(1)]
		\item Cả nhiên liệu rắn và nhiên liệu khí đều có thể tái sử dụng.
		\item Nhiên liệu rắn khi cháy sinh ra nhiều chất độc hại với môi trường hơn nhiên liệu khí.
		\item Nhiên liệu rắn và nhiên liệu khí đều cháy được và tỏa nhiều nhiệt.
		\item Nhiên liệu rắn dễ cháy hơn nhiên liệu khí.
	\end{enumerate}
	Hãy sắp xếp các phát biểu đúng thành bộ số theo số thứ tự tăng dần (ví dụ $12$, $134$, …).
	\shortans{$23$}
	\loigiai{
		\begin{itemize}
			\item $(1)$ sai vì các nhiên liệu sau khi cháy sẽ không tái sử dụng được.
			\item $(4)$ sai vì nhiên liệu khí thường dễ cháy hơn nhiên liệu rắn.
			\item Các phát biểu đúng là $(2)$ và $(3)$.
		\end{itemize}
		\textbf{\textit{Đáp số:} $23$.}
	}
\end{ex}

%%%=============SA_5=============%%%
\begin{ex}%[6K3N3-1]
	Cho các nhiên liệu: $(1)$ than đá, $(2)$ gas, $(3)$ gỗ, $(4)$ khí thiên nhiên, $(5)$ dầu hỏa, $(6)$ cồn. Hãy sắp xếp các nhiên liệu ở trạng thái khí thành bộ số theo số thứ tự tăng dần (ví dụ $12$, $134$, …).
	\shortans{$24$}
	\loigiai{
		\begin{itemize}
			\item Gas $(2)$ và khí thiên nhiên $(4)$ ở trạng thái khí.
			\item Than đá, gỗ ở trạng thái rắn; dầu hỏa, cồn ở trạng thái lỏng.
		\end{itemize}
		\textbf{\textit{Đáp số:} $24$.}
	}
\end{ex}

%%%=============SA_6=============%%%
\begin{ex}%[6K3H3-1]
	Cho các chất: $(1)$ xăng, $(2)$ cát, $(3)$ dầu diesel, $(4)$ gỗ, $(5)$ nước biển, $(6)$ khí gas. Hãy sắp xếp các chất là nhiên liệu thành bộ số theo số thứ tự tăng dần (ví dụ $12$, $134$, …).
	\shortans{$1346$}
	\loigiai{
		\begin{itemize}
			\item Xăng $(1)$, dầu diesel $(3)$, gỗ $(4)$, khí gas $(6)$ đều cháy được và tỏa nhiều nhiệt nên là nhiên liệu.
			\item Cát và nước biển không cháy được nên không phải nhiên liệu.
		\end{itemize}
		\textbf{\textit{Đáp số:} $1346$.}
	}
\end{ex}

%%%=============SA_7=============%%%
\begin{ex}%[6K3H3-1]
	Cho các nguồn nhiên liệu: $(1)$ than đá, $(2)$ dầu mỏ, $(3)$ khí thiên nhiên, $(4)$ ethanol, $(5)$ năng lượng mặt trời. Hãy sắp xếp các nhiên liệu hóa thạch thành bộ số theo số thứ tự tăng dần (ví dụ $12$, $134$, …).
	\shortans{$123$}
	\loigiai{
		\begin{itemize}
			\item Nhiên liệu hóa thạch gồm than đá $(1)$, dầu mỏ $(2)$, khí thiên nhiên $(3)$.
			\item Ethanol là nhiên liệu sinh học, năng lượng mặt trời là năng lượng tái tạo.
		\end{itemize}
		\textbf{\textit{Đáp số:} $123$.}
	}
\end{ex}

%%%=============SA_8=============%%%
\begin{ex}%[6K3V3-3]
	Cho các việc làm sau:
	\begin{enumerate}[(1)]
		\item Khóa van bình gas sau khi sử dụng xong.
		\item Đun bếp than trong phòng kín để giữ nhiệt.
		\item Chẻ nhỏ củi trước khi đun nấu.
		\item Để ngọn lửa thật to trong suốt quá trình nấu.
		\item Vặn nhỏ lửa khi thức ăn gần chín.
	\end{enumerate}
	Hãy sắp xếp các việc làm giúp sử dụng nhiên liệu an toàn, tiết kiệm thành bộ số theo số thứ tự tăng dần (ví dụ $12$, $134$, …).
	\shortans{$135$}
	\loigiai{
		\begin{itemize}
			\item $(2)$ sai vì đun bếp than trong phòng kín sinh ra khí độc carbon monoxide, gây ngạt.
			\item $(4)$ sai vì để lửa quá to gây lãng phí nhiên liệu và mất an toàn.
			\item Các việc làm đúng là $(1)$, $(3)$, $(5)$.
		\end{itemize}
		\textbf{\textit{Đáp số:} $135$.}
	}
\end{ex}

%%%=============SA_9=============%%%
\begin{ex}%[6K3H3-6]
	Cho các nguồn năng lượng: $(1)$ năng lượng mặt trời, $(2)$ than đá, $(3)$ năng lượng gió, $(4)$ dầu mỏ, $(5)$ nhiên liệu sinh học. Hãy sắp xếp các nguồn năng lượng tái tạo thành bộ số theo số thứ tự tăng dần (ví dụ $12$, $134$, …).
	\shortans{$135$}
	\loigiai{
		\begin{itemize}
			\item Nguồn năng lượng tái tạo gồm năng lượng mặt trời $(1)$, năng lượng gió $(3)$, nhiên liệu sinh học $(5)$.
			\item Than đá và dầu mỏ là nhiên liệu hóa thạch, không tái tạo.
		\end{itemize}
		\textbf{\textit{Đáp số:} $135$.}
	}
\end{ex}

%%%=============SA_10=============%%%
\begin{ex}%[6K3V3-3]
	Cho các phát biểu sau:
	\begin{enumerate}[(1)]
		\item Nhiên liệu là những chất cháy được và khi cháy tỏa nhiều nhiệt.
		\item Muốn nhiên liệu cháy hoàn toàn cần cung cấp thiếu oxygen.
		\item Tăng diện tích tiếp xúc giữa nhiên liệu và không khí giúp nhiên liệu cháy tốt hơn.
		\item An ninh năng lượng ưu tiên các nguồn năng lượng sạch và giá thành rẻ.
		\item Nên tận dụng triệt để nhiên liệu hóa thạch có sẵn trong tự nhiên.
	\end{enumerate}
	Hãy sắp xếp các phát biểu đúng thành bộ số theo số thứ tự tăng dần (ví dụ $12$, $134$, …).
	\shortans{$134$}
	\loigiai{
		\begin{itemize}
			\item $(2)$ sai vì cần cung cấp đủ (vừa đủ) oxygen thì nhiên liệu mới cháy hoàn toàn.
			\item $(5)$ sai vì tận dụng triệt để nhiên liệu hóa thạch làm cạn kiệt tài nguyên và gây ô nhiễm môi trường.
			\item Các phát biểu đúng là $(1)$, $(3)$, $(4)$.
		\end{itemize}
		\textbf{\textit{Đáp số:} $134$.}
	}
\end{ex}
\Closesolutionfile{ansexh}
\Closesolutionfile{ansex}
%\bangdapanSA{AnsSA-C03_B08_MOT_SO_NHIEN_LIEU}
\end{document}
